% ============================================================
% ShuWei Cup English paper template (shuwei_en), v1.0
% Shell: official swmcmthesis (2025 autumn). Skeleton: this system's
% problem-organized modeling paper (sensitivity / validation kept).
% Compile: copy swmcmthesis.cls + swmcmthesis.cfg next to paper.tex, then
%   pdflatex paper.tex
%   pdflatex paper.tex
% Do NOT use ctexart. Do NOT delete \label{abstract:end} / \label{appendix:start}.
% Page cap: Summary + Content + body <= 25; appendix unlimited.
% Identity: only \baominghao / \zubie; never \schoolname / \membera / advisor.
% ============================================================

\documentclass{swmcmthesis}

\usepackage{float}
\usepackage{fvextra}
\usepackage{verbatim}
\setlength{\headheight}{20pt}

\fvset{frame=single,numbers=left,numbersep=6pt,fontsize=\small,
       baselinestretch=1.05,tabsize=4,breaklines=true,breakanywhere=true}

\title{[Concise, Specific Paper Title]}
\baominghao{BSHUWEICUP26XXXXXX}  % official team control number, unchanged
\zubie{B}                         % problem letter A/B/C/D
\date{2026}                       % contest year on the Summary Sheet (not an author date)

\abstract{%
Paragraph 1 (1--3 sentences): problem setting and overall approach.

For Problem 1, we formulated [model name], solved it with [solver/algorithm],
and obtained [quantitative result: objective, gap, or policy parameters].

For Problem 2, we formulated [model name], solved it with [tool],
and obtained [quantitative result].

For Problem 3, we formulated [model name], solved it with [tool],
and obtained [quantitative result].

Keep this Summary on a single page. No school, advisor, or member names.
\label{abstract:end}%
}

\keywords{keyword1, keyword2, keyword3, keyword4}

\begin{document}

\maketitle
\tableofcontents
\newpage

% ============================================================
\section{Introduction}

\subsection{Background}

[Problem background. If the statement defines proper nouns or attachments,
a three-line table of terms may go here.]

\subsection{Problem Restatement}

\textbf{Problem 1.} [Restate as a precise mathematical task; do not copy the
statement verbatim.]

\textbf{Problem 2.} [Same.]

\textbf{Problem 3.} [Same.]

% ============================================================
\section{Problem Analysis}

% Subsection titles must be research-oriented (e.g.\ ``Coupling between
% indicator X and parameter Y''). Do not use ``Analysis of Question 1''.

\subsection{[Research-oriented title for Problem 1]}

[Core difficulty, candidate method and why, link to other problems.
About 200--400 words.]

% ============================================================
\section{Model Assumptions}

% 3--6 items, each a pure assertion. Reasons belong in Problem Analysis
% or Sensitivity, not after the item.
\begin{enumerate}
    \item Assumption (pure assertion);
    \item Assumption;
    \item Assumption.
\end{enumerate}

% ============================================================
\section{Notation}

% Single table aligned with 术语表.json. Gloss in English at first use
% in the body: ``travel cost of unit $X$ ($c_X$)''.
\begin{table}[H]
\centering
\caption{Notation}
\label{tab:notation}
\begin{tabular}{ccc}
\toprule
Symbol & Description & Unit \\
\midrule
$x_1$ & [specific meaning] & [unit, or /] \\
$x_2$ & [meaning] & [unit] \\
\bottomrule
\end{tabular}
\end{table}

% ============================================================
\section{Model Formulation and Solution}

% One problem at a time: formulate $\to$ solve $\to$ verify $\to$ interpret.
% Derivations stay in the body. Appendix holds code, file lists, and
% algebra that truly does not fit.

\subsection{[Research-oriented title for Problem 1]}

\subsubsection{Data exploration or problem structure}

[If data attachments are used: scatter/distribution $\to$ pattern $\to$
outliers and preprocessing. Otherwise: structural reading of the instance.]

\subsubsection{Model formulation}

[How assumptions become mathematics. State the form, then define each term.]
\begin{equation}
y = f(x_1, x_2, \ldots; \theta)
\end{equation}
where $\theta$ is [meaning] and $x_i$ is [meaning].

\subsubsection{Solution method}

[At least two algorithm candidates and the selection reason ($\ge 1$ paragraph).]

\begin{table}[H]
\centering
\caption{Model / algorithm comparison}
\label{tab:model_compare}
\begin{tabular}{lcccc}
\toprule
Model / algorithm & Objective & Optimality & Time & Note \\
\midrule
MIP exact & {value} & MIPGap = {x}\% & {s} & small / exact \\
Heuristic & {value} & gap to exact {x}\% & {s} & large / real-time \\
\bottomrule
\end{tabular}
\end{table}

\subsubsection{Verification and result analysis}

[Optimality: MIPGap / exhaustive enumeration / heuristic statistics.
Model validation, when randomness or a steady-state assumption is present:
Monte Carlo replications, convergence, deviation from the analytic value,
and a cause if the deviation exceeds 5\%. The two checks are not substitutes.]

[Table + figure + $\ge 150$ words: trend / peak / kink, link to the model,
and practical meaning.]

\begin{figure}[H]
\centering
\includegraphics[width=0.9\textwidth]{fig_result_p1.pdf}
\caption{Problem 1 main result}
\label{fig:result_p1}
\end{figure}

% Repeat the same subsection pattern for remaining problems.
\subsection{[Research-oriented title for Problem 2]}

[Formulate, solve, verify, interpret. Cite Fig.~\ref{fig:result_p2}.]

\begin{figure}[H]
\centering
\includegraphics[width=0.9\textwidth]{fig_result_p2.pdf}
\caption{Problem 2 main result}
\label{fig:result_p2}
\end{figure}

\subsection{[Research-oriented title for Problem 3]}

[Formulate, solve, verify, interpret. Cite Fig.~\ref{fig:result_p3}.]

\begin{figure}[H]
\centering
\includegraphics[width=0.9\textwidth]{fig_result_p3.pdf}
\caption{Problem 3 main result}
\label{fig:result_p3}
\end{figure}

% ============================================================
\section{Sensitivity and Error Analysis}

\subsection{Sensitivity analysis}

[Perturb core parameters by $\pm 5\%/\pm 10\%/\pm 20\%$. Tornado chart plus
ranking. Ranges must match the code.]

\begin{figure}[H]
\centering
\includegraphics[width=0.9\textwidth]{fig_sensitivity.pdf}
\caption{Parameter sensitivity tornado chart}
\label{fig:tornado}
\end{figure}

\begin{table}[H]
\centering
\caption{Sensitivity summary}
\label{tab:sensitivity}
\begin{tabular}{ccccc}
\toprule
Parameter & Base & Perturbation & Change in result & Rank \\
\midrule
{param 1} & {base} & $\pm 5\%$ & $+{x}\%$ & 1 (most sensitive) \\
{param 2} & {base} & $\pm 5\%$ & $-{x}\%$ & 2 \\
\bottomrule
\end{tabular}
\end{table}

\begin{figure}[H]
\centering
\begin{minipage}{0.48\textwidth}
\centering
\includegraphics[width=\textwidth]{fig_a.pdf}\\
{\small (a) Case A}
\end{minipage}
\hspace{0.04\textwidth}
\begin{minipage}{0.48\textwidth}
\centering
\includegraphics[width=\textwidth]{fig_b.pdf}\\
{\small (b) Case B}
\end{minipage}
\caption{Paired comparison (subfigures (a)(b); caption states the claim)}
\label{fig:twin}
\end{figure}

\subsection{Policy switching thresholds (optional)}

[If the optimal policy switches with a parameter, report the critical value
with a figure or table sourced from the solver.]

\subsection{Active constraints (optional)}

[At the optimum, each constraint active/slack; gain from relaxing a tight one.]

\subsection{Error sources and robustness}

[Source, magnitude, impact, acceptability; robustness check.]

\begin{figure}[H]
\centering
\includegraphics[width=0.9\textwidth]{fig_flow.pdf}
\caption{Solution workflow}
\label{fig:flow}
\end{figure}

% ============================================================
\section{Model Evaluation}

\subsection{Strengths}

[Strengths.]

\subsection{Limitations}

[Limitations. Absorb METHOD \texttt{PASS\_WITH\_NOTES} tickets here.]

\subsection{Improvements and extensions}

[Improvements and transferable settings.]

% Do not add a section that critiques the problem statement.

% ============================================================
\begin{thebibliography}{99}
% Official style (2025 autumn): numbered in citation order; books / journals /
% web as specified. Not GB/T 7714. Years must not post-date the contest.
\bibitem{ref1} Author, Book Title, Place: Publisher, 2000.
\bibitem{ref2} Author, Paper Title, Journal Name, vol.(no.): pages, 2010.
\bibitem{ref3} Author, Resource Title, URL, accessed YYYY-MM-DD.
\bibitem{ref4} Dantzig G B, Linear Programming and Extensions, Princeton: Princeton University Press, 1963.
\bibitem{ref5} Bertsimas D, Tsitsiklis J N, Introduction to Linear Optimization, Belmont: Athena Scientific, 1997.
\bibitem{ref6} Boyd S, Vandenberghe L, Convex Optimization, Cambridge: Cambridge University Press, 2004.
\bibitem{ref7} Wolsey L A, Integer Programming, New York: Wiley, 1998.
\bibitem{ref8} Ahuja R K, Magnanti T L, Orlin J B, Network Flows, Englewood Cliffs: Prentice Hall, 1993.
\bibitem{ref9} Lawler E L, Combinatorial Optimization: Networks and Matroids, New York: Holt, Rinehart and Winston, 1976.
\bibitem{ref10} Nemhauser G L, Wolsey L A, Integer and Combinatorial Optimization, New York: Wiley, 1988.
\end{thebibliography}

% Body page count: appendix:start - 1 (Summary + Content + body). Cap 25.
% Do not rename the label.
\clearpage
\appendix
\section{Supporting Materials}
\label{appendix:start}

\begin{table}[H]
\centering
\caption{Supporting materials}
\label{tab:support}
\begin{tabular}{p{5cm}p{9cm}}
\toprule
\multicolumn{1}{c}{File} & \multicolumn{1}{c}{Description} \\
\midrule
Code & [n] Python files (entry: main.py / problem1\_solver.py, \ldots) \\
Results & [CSV outputs + sensitivity + figures] \\
{other.pdf} & [e.g.\ AI-tool disclosure] \\
\bottomrule
\end{tabular}
\end{table}

\section{Source Code}

[All runnable source used in the paper. \texttt{\textbackslash VerbatimInput}
must point at \texttt{project/} files; do not rewrite the listing by hand.]

\VerbatimInput{project/problem1_solver.py}

\section{Supplementary derivations}

[Only intermediate algebra that cannot fit in the body. The body must cite
``see Appendix C''. Do not write ``the derivation is omitted''.]

\end{document}
